\subsection{Cell-selective expression: cohort and measurement}
\label{app:cell-selective}
The Siraj HepG2/K562 saturation assay \citep{siraj2026functional} supplies
167,958 substitutions, 280 parental backgrounds and 127 loci. One locus with
four indel-containing backgrounds lacks prepared physical-count coverage, excluding 2,334
substitutions. No activity, agreement or model-score filter is added.
Each parent and requested target cell defines an opportunity. Both cells serve
as target, and all measured substitutions, including the natural-allele
substitution, compete with no edit. Unlike response matching, this task permits
direct reversions.

For sequence $s$, cell $c$ and repeat $r$, define
\begin{equation}
 z_{s,c,r}=\log_2\frac{\mathrm{RNA}_{s,c,r}+0.5}{\mathrm{DNA}_{s,r}+0.5},
 \qquad
 \Delta_c(s)=\frac{1}{5}\sum_{r=1}^{5}(z_{s,c,r}-z_{m0,c,r}),
\end{equation}
where $m0$ is the physical unmutated parent. Effects average log ratios against
that parent; counts are not pooled, and no saturation-library baseline is used.
DNA counts are shared between cells.
Subsets A (repeats 1--2) and B (3--5) recompute their own parent contrasts and
test robustness within one library, not independent studies.

\paragraph[Predicting effects.]{\phantomsection\label{app:cell-selective-fitting}Predicting effects.}
Figure~\ref{fig:control}C compares AlphaGenome and Enformer accessibility,
NTv3 100M/650M post readouts, 3--6-mer ridge and the task-trained one-hot ResNet.
AlphaGenome and Enformer effects are aggregated over 501-bp and 1,024-bp output windows, respectively. To calibrate native
accessibility, each cell and fold receives a nonnegative through-zero slope:
\begin{equation}
 b=\max\!\left(0,\frac{\sum_i w_i a_i\Delta_i}{\sum_i w_i a_i^2}\right).
\end{equation}
Here $a_i$ is the raw native effect; a zero denominator gives slope zero.
Five deterministic folds keep complete loci together.
Fitting excludes evaluation-fold loci and weights training loci equally, then
backgrounds within loci and candidates within backgrounds equally. These native outputs are
assay-calibrated, not zero-shot.

Fitted predictions combine the predicted mean effect across cells with the predicted cell
contrast: HepG2 is the mean plus half the HepG2-minus-K562 contrast, and K562
is the mean minus half that contrast. Readouts fitted to the original
assay estimator are evaluated against physical-parent effects without recalibration.
Supervision, context and readout training differ, so comparisons concern these
pipelines rather than architecture alone.

\paragraph{Choosing and evaluating edits.}
For fixed $\lambda\in\{0,0.25,0.5,1,2,4\}$, nominate the substitution maximizing
$\widehat\Delta_{\rm target}-\lambda|\widehat\Delta_{\rm other}|$.
Abstain at maxima of at most $10^{-10}$; otherwise share weight equally among
ties within absolute tolerance $10^{-10}$ and zero relative tolerance.
The primary setting is $\lambda=1$, chosen without measured outcomes.
Evaluate target effect $T$, absolute other-cell effect $O$ and $U=T-O$.
Take absolute values before averaging ties, backgrounds or cells. Average
opportunities within loci, then the 127 loci equally. No edit contributes zero;
uniform choice uses the exact expectation over eligible substitutions.

\paragraph{Prediction-matched strength control.}
Weaker edits can reduce desired and unwanted effects together. The control
therefore uses the same attempted opportunities and matches predicted target
response within each one. It mixes six target-only policies: uniform choice
in relative predicted-strength bands $(0,.25]$, $(.25,.5]$, $(.5,.75]$, $(.75,1]$,
minimum positive strength and maximum strength. Empty bands use candidates
nearest the midpoint. Interpolation between adjacent component means matches
the selective policy's expected predicted target effect. Other-cell scores do
not rank candidates within components; measured responses do not construct the
mixture. Matching predictions does not establish equal measured effects.

\paragraph{Uncertainty and results.}
The 24-contrast family tests six models on utility versus no edit,
utility versus matched control, target-effect difference and other-cell reduction.
We use 50,000 shared locus-bootstrap draws and Bonferroni tails
$0.025/24$. Fits, slopes and nominations remain fixed; intervals omit training
and calibration-fitting variation. Mean target equivalence requires the adjusted
interval inside $\pm0.05$ log$_2$; no pipeline establishes it.
AlphaGenome, both post-trained NTv3 readouts and Enformer achieve positive utility;
AlphaGenome and both NTv3 readouts also improve on their strength controls.
Enformer's control gain is inconclusive: $0.055\;[-0.011,0.120]$.
These within-model comparisons do not establish a between-model ranking.
NTv3 650M pre attempts only 10.24\% of opportunities, explaining its near-zero utility.

\paragraph{Experimental reference.}
Select using subset A's measurements and evaluate in B, then reverse, retaining the
same candidate sets, penalties, no-edit option and strength control. Both directions must
have positive utility and gains over target-only and matched-strength controls;
averaged utility and target-only gain must also have positive lower pointwise
95\% locus-bootstrap bounds. These requirements hold: mean utility is
$0.962\;[0.829,1.092]$, target-only gain $0.371\;[0.297,0.448]$, and matched-control
gain $0.292\;[0.239,0.347]$. On the same subsets, reference-minus-AlphaGenome
utility is $0.603$, positive at 108 loci, negative at 18 and tied at one.
The reference uses measured candidate responses and shares oligos, parents and
library construction across subsets. It shows that selective edits can be found
across repeat subsets.
It is not a noise ceiling or independent replication.

\paragraph{Protocol.}
This evaluation is retrospective. The protocol, cohort and interfaces were fixed
before native-slope fitting and evaluation; nominations were fixed after calibration
and before scoring.
